% Strategic Management – Midterm Cheat Sheet
% Compile with: pdflatex midterm-cheat-sheet.tex

\documentclass[8pt,landscape,a4paper]{article}
\usepackage[utf8]{inputenc}
\usepackage[T1]{fontenc}
\usepackage{microtype}
\usepackage[left=0.8cm,right=2.6cm,top=0.8cm,bottom=0.8cm,marginparwidth=2cm,marginparsep=0.3em]{geometry}
\usepackage{multicol}
\usepackage{marginnote}
\usepackage{enumitem}
\setlist{nosep,leftmargin=*}
\usepackage{amsmath}
\usepackage{booktabs}
\usepackage{parskip}
\usepackage{xcolor}
\usepackage{soul}
\usepackage{fancyhdr}
\definecolor{lightpink}{RGB}{255,182,193}
\definecolor{lightpeach}{RGB}{255,218,185}
\definecolor{noteblue}{RGB}{173,216,230}
\newcommand{\handnote}[1]{\marginnote{\raggedright\footnotesize\color{noteblue}#1\vspace{0.4em}}}
\newcommand{\hlpink}[1]{\sethlcolor{lightpink}\hl{#1}\sethlcolor{lightpeach}}
\newcommand{\hlpeach}[1]{\sethlcolor{lightpeach}\hl{#1}}
\sethlcolor{lightpeach}
\setlength{\parskip}{0.2em}
\setlength{\columnsep}{0.5em}
\pagestyle{fancy}
\fancyhf{}
\fancyfoot[R]{\raisebox{0.8em}{\thepage}}
\fancypagestyle{plain}{\fancyhf{}\fancyfoot[R]{\raisebox{0.8em}{\thepage}}}
\raggedright

\begin{document}
\begin{center}
{\normalsize\bfseries Strategic Management -- Midterm Cheat Sheet}\\[0.2em]
\footnotesize\textit{Use for quick recall; in exam: define terms, give one example, tie to competitive advantage.}\\[0.4em]
\fbox{\parbox{0.92\textwidth}{\centering\footnotesize
\textcolor{black!80}{\textbf{Key:}} \colorbox{lightpink}{\strut\ \ \ } \textbf{Pink} = definitions \quad
\colorbox{lightpeach}{\strut\ \ \ } \textbf{Peach} = frameworks \& tips \quad
\colorbox{noteblue}{\strut\ \ \ } \textbf{Blue} = margin notes
\par}}
\end{center}
\vspace{-0.6em}

\begin{multicols}{2}

\section*{Unit 0 -- Course \& Big Ideas}\handnote{Long-term, scarce resources -- know these by heart.}
\begin{itemize}[leftmargin=*,nosep]
\item \textbf{\hlpink{Strategy}}: long-term, hard-to-reverse decisions; compete for \textbf{\hlpink{scarce resources}} (customers, talent, capital). Limited ``spots'' in the market; we compete for them.
\item \textbf{Strategic vs tactical}: Strategic = long-term, shapes future, costly to undo (what to study, where to live, career change, mergers, entering a country). ``No second chance'' for many. Tactical = short-term, day-to-day (scheduling, small price changes).
\item \textbf{Forecasting}: future uncertain (economy, politics, climate, health); we imagine \textbf{scenarios}, we don't know exact outcomes. Theory is a \textbf{tool}, not the answer.
\item \textbf{Key rules}: \textbf{\hlpink{Know yourself}} = strengths, weaknesses, goals, constraints. \textbf{\hlpink{Know your environment}} = opportunities and threats (macro + industry).
\item \textbf{Social science}: politics, power, culture, stakeholders matter. Rational analysis has limits. No single ``right'' answer.\handnote{Not just shareholders! Workers, society, env.}
\item \textbf{Stakeholders} (not just shareholders): workers, community, society, environment, state/family owners; long-term success needs more than profit. Debate: \textbf{shareholder primacy} vs \textbf{stakeholder/CSR/sustainability}.
\end{itemize}

\section*{Unit 1 -- Strategy \& The Art of War}\handnote{Apply: terrain? enemy? what does firm know?}
\begin{itemize}[leftmargin=*,nosep]
\item \textbf{Strategos} (Greek) = the general; art of commanding the army. Linked to war \& competition for limited positions and resources.
\item \textbf{Why ``war''?} On a battlefield, armies compete for territory; in business, firms compete for customers, talent, capital. Same logic: limited spots, someone wins, someone loses.
\item \textbf{Strategic decisions in life}: what to study, career, where to live, partner, children $\Rightarrow$ big effects, no easy second chance. Think 10+ years ahead.
\item \textbf{Sun Tzu analogies}: Battlefield $\to$ \textbf{market} (where competition happens). Enemy $\to$ \textbf{competitors/obstacles}. General $\to$ \textbf{leader/strategist}. Terrain $\to$ \textbf{industry structure, rules, geography}.
\item \textbf{Core lesson}: If you \hlpeach{know yourself} (resources, capabilities, goals), \hlpeach{know the enemy} (competitors, their moves), and \hlpeach{know the terrain} (market, industry, regulation), your chances of success are much higher.
\item \textit{\hlpeach{Exam tip}}: In a case, ask: What is the terrain? Who are the enemies? What does the firm know or not know?
\end{itemize}

\section*{Unit 2 -- Objectives, Values, Performance, SWOT}\handnote{Mission = who. Vision = where. Objectives = how.}

\subsection*{A. Strategy as intention}
Strategy = where we want to go + how we get there. Long-term = often $>$10 years.\handnote{Always compare performance to something!}

\subsection*{B. Mission, vision, objectives}
\begin{itemize}[leftmargin=*,nosep]
\item \textbf{\hlpink{Mission}} = Who am I? What do I do? \textbf{Current state.} Products/services + target market(s). One sentence (e.g.\ ``I produce and sell cars'').
\item \textbf{\hlpink{Vision}} = What do I want to become? \textbf{Future state}, long-term (e.g.\ ``Leader in electric vehicles in Europe by 2035'').
\item \textbf{\hlpink{Strategic objectives}} = Concrete steps from mission to vision (milestones, targets).
\item \textbf{Objectives} (general): profit, growth, survival, reputation. \textbf{Values}: what firm will/won't do (ethics, CSR). Constrain \textit{how} you compete.
\item \textbf{Performance}: always ``compared to what?'' \textbf{Referent group} = similar firms in industry. \textbf{History} = firm's own past. \textbf{Intention} = chosen constraints (e.g.\ closing Sundays).
\end{itemize}

\subsection*{C. Measuring performance}
\begin{itemize}[leftmargin=*,nosep]
\item \textbf{Accounting} (non-listed/family): \textbf{\hlpink{ROA}} = profit/total assets (efficiency); \textbf{\hlpink{ROE}} = profit/equity (return for shareholders). A number only makes sense vs industry average, main rival, or own past. No absolute ``good'' number.
\item \textbf{Market} (listed): Shareholder return = (price end $-$ price start) + dividends. \textbf{Opportunity cost} = next best alternative (other firms, index, risk-free). Satisfied when return $\ge$ opportunity cost.
\item \textbf{\hlpink{Cost of equity}} = minimum return shareholders require: $\text{Cost of Equity} = R_f + \beta \times (\text{Market Return} - R_f)$. \textbf{$R_f$} = risk-free rate (e.g.\ govt bonds). \textbf{$\beta$} = sensitivity to market. \textbf{Market Return} = expected market return (e.g.\ index). \hlpink{Satisfaction}: return $>$ cost of equity (or better than industry average).
\end{itemize}\handnote{Formula: Rf + beta (Market - Rf). Compare vs industry.}

\subsection*{D--E. Competitive advantage \& Porter}
\begin{itemize}[leftmargin=*,nosep]
\item \textbf{\hlpink{CA}} at each \textbf{transaction} when customer chooses you instead of a rival or not buying. Built sale by sale; over many transactions it looks ``sustained''.
\item \textbf{Low-cost}: win on price; customers price-sensitive (elastic). E.g.\ budget airlines, generic products.
\item \textbf{Differentiation}: customers strongly prefer your offer (brand, quality, design, service); inelastic. E.g.\ luxury brands, Apple.
\item \textbf{Focus}: narrow segment (niche); can combine with low-cost or differentiation. \textbf{Stuck in the middle}: Porter warned---neither low-cost nor differentiated $\Rightarrow$ worst of both (low margins, no loyalty).
\end{itemize}\handnote{Low-cost, diff, focus. Never stuck in the middle!}

\subsection*{F. Governance \& CSR}
Companies = social organisations. Stakeholders: shareholders (return), managers (power, pay), workers (wages, conditions), customers (quality, price), society (jobs, taxes, pollution). \textbf{Agency problem}: managers $\neq$ owners; need carrots (stock options, bonuses) and sticks (board oversight). \textbf{Externalities}: negative (e.g.\ pollution) and positive (jobs, taxes) often not fully priced. Debate: shareholder primacy vs stakeholder/CSR/sustainability (climate, social expectations).\handnote{Agency: managers vs owners. Externalities = not priced.}

\subsection*{G. SWOT done right}
S/W = \textbf{internal} (resources, capabilities). O/T = \textbf{external} (environment, industry, PESTEL). \textbf{Bad}: vague (``our people are our strength'', ``the market is growing'')---no link to why customers choose you. \textbf{Good}: specific strengths/weaknesses tied to CA (e.g.\ ``logistics allows 24h delivery'', ``brand in segment X'', ``location next to university''); concrete O/T (e.g.\ ``new regulation favours renewables'', ``two new entrants 2024''). \textit{\hlpeach{Exam tip}}: What do we have (S/W)? What's outside (O/T)? How do S and O combine?\handnote{S/W = internal. O/T = external. Be specific!}

\section*{Unit 3 -- PESTEL, Five Forces, Clusters}\handnote{Macro + industry. Can't control but must adapt.}

\subsection*{Environment}
External factors firm \textbf{cannot control} but that \textbf{influence} behaviour and results. Managers' job: detect opportunities and threats. Two levels: \textbf{macro} (country/world) and \textbf{industry} (specific arena).

\subsection*{\hlpeach{PESTEL}}
\begin{itemize}[leftmargin=*,nosep]
\item \textbf{P}olitical: govt, stability, policy, taxes, trade, war/peace. E.g.\ change of government $\to$ new subsidies or tariffs.
\item \textbf{E}conomic: growth, inflation, interest rates, unemployment, exchange rates. E.g.\ recession $\to$ less demand; strong euro $\to$ export harder.
\item \textbf{S}ocial: demographics, culture, lifestyles, education, inequality. E.g.\ ageing population $\to$ health demand; remote work $\to$ tech demand.
\item \textbf{T}echnological: R\&D, new tech, digitalisation. E.g.\ AI, automation $\to$ new products, displaced jobs.
\item \textbf{E}cological: climate, resources, pollution. E.g.\ regulation on emissions $\to$ cost or green tech opportunities.
\item \textbf{L}egal: labour, consumer, competition, environment, IP. E.g.\ GDPR, minimum wage.
\end{itemize}
PESTEL changes $\Rightarrow$ new opportunities \& threats $\Rightarrow$ firms must adapt.\handnote{Six letters: P E S T E L. One example each.}

\subsection*{\hlpeach{Porter's Five Forces}}
Aim: \textbf{industry attractiveness} (profit potential). Strong forces = less attractive; weak = more attractive. Always from \textbf{firm's perspective} (good for buyers/suppliers = bad for firm).
\begin{enumerate}[leftmargin=*,nosep]
\item \textbf{Rivalry}: high when many competitors, low growth, commoditised (low differentiation), high exit barriers, high fixed costs. Result: price wars, lower margins.
\item \textbf{New entrants}: easy entry $\Rightarrow$ more firms $\Rightarrow$ more rivalry. \textbf{Barriers}: natural (economies of scale, know-how, tech, learning curve); artificial (licenses, regulation, tariffs, patents); capital requirements; access to distribution; strong brands; expected retaliation.
\item \textbf{Substitutes}: products/tech that satisfy same need (e.g.\ train vs plane; streaming vs cinema). Disruptive technology or consumer behaviour can increase this force.
\item \textbf{Buyer power}: high when few large buyers (concentration); low switching cost; they have full info on your costs; they can backward integrate (make it themselves). Result: they push down price $\Rightarrow$ less attractive for you.
\item \textbf{Supplier power}: high when few suppliers; you depend on their specific input (differentiation); they have info on your business; they can forward integrate (become your competitor). Result: they push up input prices $\Rightarrow$ less attractive.
\end{enumerate}
Goal: \textbf{\hlpeach{no organising power}} for either buyers or suppliers (keep margin and control).\handnote{Strong force = bad for you. Weak = good.}

\subsection*{Limits \& extras}
Performance = \textbf{\hlpeach{industry effect}} (attractiveness) + \textbf{\hlpeach{firm effect}} (own strategy and resources). Firm can do well in unattractive industry (strong resources) or badly in attractive one (weak strategy). Not all five forces equally relevant; focus on strongest 2--3. \textbf{Complements}: products that go with yours (e.g.\ software + hardware) can increase demand. \textbf{Sixth force}: public admin/regulators (Europe: energy, telecoms, health). \textbf{Segments \& strategic groups}: not all players face same forces. Five Forces = simplification of reality.\handnote{Industry + firm effect. Not all forces equal -- pick 2--3.}

\subsection*{Clusters}\handnote{Geography! Firms + suppliers + unis. Silicon Valley, Toyota city.}
\textbf{Geographical agglomeration} of firms in one activity + suppliers + related industries + universities/research + support (e.g.\ Silicon Valley, Toyota city, Italian textile districts). Benefits: talent attraction, entrepreneurship, technological innovation, access to finance and networks. Proximity reduces transaction costs and speeds learning.

\subsection*{Industry definition \& attractiveness}
Define from \textbf{customer's point of view}: not just ``all car makers'' but all offers that \textbf{meet the same need/wish} (including substitutes, luxury experiences). \textbf{Attractive} when opportunities $>$ threats $\Rightarrow$ easier profit, growth, survival. \textbf{Unattractive} when threats $>$ opportunities. Manager's job: enter attractive industries, avoid/exit bad ones when possible.\handnote{Customer need = industry. Attractive = O $>$ T.}

\subsection*{Climate (PESTEL case)}\handnote{Arctic example: all 6 PESTEL.}
Climate change = economic and strategic, not only environmental. Melting Arctic/Greenland: new sea routes, trade patterns; resources (oil, gas, rare earths) more accessible; geopolitical competition, legal disputes. PESTEL: Political (sovereignty, military), Economic (routes, costs), Social (local communities), Technological (ice-breaking, extraction), Environmental (fragile ecosystem), Legal (UNCLOS, maritime law).

\section*{Unit 4 -- \hlpeach{RBV} \& Internal Resources}\handnote{Internal view! 70-80\% = you, not industry.}
\begin{itemize}[leftmargin=*,nosep]
\item \textbf{RBV} (Barney and others): strategy from \textbf{what firm has inside} (resources \& capabilities), not only external (PESTEL, Five Forces). Internal view \textbf{complements} external view.
\item \hlpeach{70--80\%} of success/failure = \textbf{internal} factors (management, people, resources, capabilities), not just industry or macro.
\item You don't need to be ``the best''; need to be \textbf{better than competitors} and \textbf{sustain} that advantage over time. \textbf{\hlpink{Sustainable CA}} = advantage that lasts because rivals cannot easily copy or substitute it.
\item \textbf{Resources} = what you \textit{have} (assets, people, brand, patents). \textbf{Capabilities} = what you \textit{do} with them (processes, routines, coordination). Both matter for RBV.
\end{itemize}\handnote{Have = resources. Do = capabilities.}

\subsection*{\hlpeach{Value Chain} (Porter, 1985)}
Primary: logistics, operations, marketing, service. Supporting: procurement, HR, tech. Advantage from activities, horizontal links (inside firm), vertical links (suppliers/buyers). E.g.\ Mercadona, Ford Valencia.

\subsection*{\hlpeach{VRIN / VRIO}}\handnote{All 4 + Org needed for sustained CA.}
\textbf{V}aluable $\Rightarrow$ helps exploit opportunities or neutralise threats (creates value). \textbf{R}are $\Rightarrow$ few or no competitors have it. \textbf{I}nimitable $\Rightarrow$ difficult or costly to copy (e.g.\ unique culture, tacit knowledge). \textbf{N}on-substitutable $\Rightarrow$ cannot be replaced by something else that does the same job. \textbf{O} (VRIO) = \textbf{Organisation}: structure and processes to \textbf{exploit} the resource (otherwise V+R+I+N wasted). Only resources that meet all four (V+R+I+N) and are well organised (O) support \textbf{sustainable} CA. If valuable but not rare, or rare but easy to copy, advantage is temporary.

\subsection*{Tangible vs.\ intangible}
\begin{center}
\footnotesize
\begin{tabular}{@{}lll@{}}
\toprule
\textbf{Aspect} & \textbf{Tangible} & \textbf{Intangible} \\
\midrule
What & Physical, financial & Human, tech, reputation, brand \\
Where & In financial statements & Often not in books \\
Control & Easier & Harder; ``double loss'' when people leave \\
\bottomrule
\end{tabular}
\end{center}
\hlpeach{$\sim$70\%} of value can be intangibles.\handnote{Intangibles often not in books!}

\subsection*{Human-based resources}
Knowledge, motivation, behaviour, loyalty ``in the person''; firm does \textbf{not own} the person. They can leave. When key people leave $\Rightarrow$ \textbf{\hlpink{double loss}}\handnote{Person + knowledge gone!} (you lose both the person and the knowledge; what was in their heads is gone). \textbf{Retention}: culture, identity, belonging, incentives (e.g.\ ``we are champions''), not only pay. If resource is \textbf{complementary} to others and hard to imitate, competitors cannot easily ``buy'' the same advantage (e.g.\ star player + team + culture + system). \textbf{Example---football}: Real Madrid, Barça, Man City; players = human resources; Ronaldo, Messi = rare, VRIO. Reputation (club brand) goes up or down with results and behaviour. Most value \textbf{intangible} (not on balance sheet).

\subsection*{Reputation / brand}\handnote{Brand up or down with behaviour.}
Use of brand name can \textbf{increase or decrease} value over time (quality, scandals, consistency). Intangible; not fully reflected in financial statements. Same for human capital and many capabilities.

\section*{Case 1 -- Madonna \& Lola Flores}\handnote{Economic vs social success. Both = human resources.}

\subsection*{Economic vs social success}
\textbf{Madonna}: one of the most profitable artists; \textbf{economic} success (revenue, profitability, global reach). \textbf{Lola Flores}: huge \textbf{social and cultural} contribution---working class, radio, fashion, identity; ``did more than profit'' for society. Both = \hlpink{human-based strategic resources}; we can assess value and rareness in economic \& social terms.

\subsection*{\hlpeach{Grant's four elements} (apply to any case)}\handnote{Goals, env, resources VRIO?, implementation.}
(1) \textbf{Simple, consistent, long-term goals}. (2) \textbf{Profound understanding of the environment} (industry, culture). (3) \textbf{Objective appraisal of resources and capabilities}. (4) \textbf{Effective implementation} (organisation, choices, actions).

\subsection*{Madonna}
\textbf{Goals}: ``Queen of Pop'', global cultural icon. \textbf{Environment}: New York; talent not enough; industry intensely competitive, volatile; need differentiation and reinvention. \textbf{Resources}: discipline, charisma, work ethic; reinvention of personas/styles; strong personal brand; network (choreographers, designers, producers); instinct for trends. \textbf{Implementation}: autonomy over artistic direction; controversy \& media to stay visible; teams for each phase. \textbf{VRIO examples}: reinvention ability (valuable, rare, hard to copy); global brand; creative network; instinct for trends + discipline.\handnote{Reinvention = VRIO. Differentiation.}

\subsection*{Lola Flores}
\textbf{Goals}: succeed from modest beginnings; central Spanish cultural figure. \textbf{Environment}: economic hardship, limited roles for women; differentiated by blending flamenco, copla, popular theatre. \textbf{Resources}: charisma, stage presence; ``La Faraona'' identity (gipsy roots); hybrid style; family dynasty (``Flores Clan''). \textbf{Implementation}: film, radio, TV; media presence, humour, authenticity; family in performances. \textbf{VRIO examples}: charisma \& emotional connection; ``La Faraona'' identity; fusion style (flamenco + copla + theatre); family clan amplifying brand.\handnote{Identity + family = rare, hard to copy.}

\subsection*{\hlpeach{Exam tip}}\handnote{Goals $\to$ env $\to$ VRIO $\to$ implementation.}
Apply Grant's four to any case: (1) What are the goals? (2) How well does the firm understand the environment? (3) What resources/capabilities (VRIO)? (4) How effective is implementation?

\section*{Quick revision -- Can you explain\ldots?}\handnote{Run through these before the exam!}
\begin{itemize}[leftmargin=*,nosep]
\item \textbf{Strategy} in one sentence (long-term, scarce resources, competition).
\item \textbf{Strategic vs tactical} with one example each.
\item \textbf{Art of War}: know yourself, enemy, terrain $\Rightarrow$ apply to a firm.
\item \textbf{Mission vs vision vs objectives} (current state, future state, steps).
\item \textbf{ROA vs ROE} and when to use accounting vs market valuation.
\item \textbf{Cost of equity} formula and what ``shareholders satisfied'' means.
\item \textbf{CA} at the transaction level (customer chooses you).
\item \textbf{Porter's three strategies} (low-cost, differentiation, focus) with examples.
\item \textbf{SWOT done right} (specific, linked to why customers choose you).
\item \textbf{PESTEL}---six factors and one example each.
\item \textbf{Five Forces}---all five; what makes each force strong (bad for you)?
\item \textbf{Barriers to entry}: natural vs artificial.
\item \textbf{Buyer and supplier power}: concentration, differentiation, information, vertical integration.
\item \textbf{Industry attractiveness} and ``industry effect + firm effect''.
\item \textbf{RBV}: why internal resources matter (70--80\%); \textbf{VRIN/VRIO}.
\item \textbf{Tangible vs intangible} (what, where in statements, control).
\item \textbf{Human resources}: why ``double loss'' when they leave; retention.
\item \textbf{Madonna vs Lola Flores}: economic vs social success; both as VRIO human resources.
\end{itemize}\handnote{If you can explain all these, you're ready.}

\section*{\hlpeach{Common question angles}}\handnote{Memorise these 4 question types!}
\begin{itemize}[leftmargin=*,nosep]
\item \textbf{``Compare internal and external analysis.''} External = PESTEL + Five Forces (opportunities, threats, industry attractiveness). Internal = RBV, resources, capabilities, VRIO (strengths, weaknesses). Both needed for strategy.
\item \textbf{``Why might a firm have CA in an unattractive industry?''} Strong \textbf{internal} resources and capabilities (VRIO) can outweigh a bad industry. E.g.\ niche player with unique know-how.
\item \textbf{``Why do financial statements underestimate firm value?''} Many key resources are \textbf{intangible} (human capital, reputation, culture) and not fully reflected; $\sim$70\% of value can sit in intangibles.
\item \textbf{``Apply Grant's four elements to [case].''} Goals $\to$ environment understanding $\to$ resources/capabilities (VRIO?) $\to$ implementation.
\end{itemize}

\end{multicols}

\newpage
\vspace*{0.3em}
\noindent\color{noteblue}
\setlength{\parskip}{0.3em}
\begin{multicols}{2}

\textbf{Unit 0 -- Course \& Big Ideas}\\
Long-term, scarce resources -- know these by heart.

Not just shareholders! Workers, society, env.

\textbf{Unit 1 -- Strategy \& The Art of War}\\
Apply: terrain? enemy? what does firm know?

\textbf{Unit 2 -- Objectives, Values, Performance, SWOT}\\
Mission = who. Vision = where. Objectives = how.

Always compare performance to something!

Formula: Rf + beta (Market - Rf). Compare vs industry.

Low-cost, diff, focus. Never stuck in the middle!

Agency: managers vs owners. Externalities = not priced.

S/W = internal. O/T = external. Be specific!

\textbf{Unit 3 -- PESTEL, Five Forces, Clusters}\\
Macro + industry. Can't control but must adapt.

Six letters: P E S T E L. One example each.

Strong force = bad for you. Weak = good.

Industry + firm effect. Not all forces equal -- pick 2--3.

Geography! Firms + suppliers + unis. Silicon Valley, Toyota city.

Customer need = industry. Attractive = O $>$ T.

Arctic example: all 6 PESTEL.

\textbf{Unit 4 -- RBV \& Internal Resources}\\
70-80\% = you. Value chain: primary + supporting; horizontal + vertical links. VRIO. Intangibles.

Have = resources. Do = capabilities.

All 4 + Org needed for sustained CA.

Intangibles often not in books!

Person + knowledge gone!

Brand up or down with behaviour.

\textbf{Case 1 -- Madonna \& Lola Flores}\\
Economic vs social success. Both = human resources.

Goals, env, resources VRIO?, implementation.

Reinvention = VRIO. Differentiation.

Identity + family = rare, hard to copy.

Goals $\to$ env $\to$ VRIO $\to$ implementation.

\textbf{Quick revision}\\
Run through these before the exam!

If you can explain all these, you're ready.

\textbf{Common question angles}\\
Memorise these 4 question types!

\end{multicols}
\vfill
\normalsize\normalfont\color{black}
\end{document}
